\section{Quantum Dynamics}
\subsubsection*{Funzioni di un'operatore}
Per proseguire con la lettura di questo capitolo è necessario definire il comportamento di una funzione di operatori. Data una funzione analitica e un operatore $\A$, l'operatore $f(\A)$ è definito tramite espansione in serie di Taylor:
\[
    f(\hat{A}) = \sum_{n=0}^{\infty} \frac{f^{(n)}(x_0)}{n!} \A^n
\]
Per esempio, se si pone $f(\A) = e^{\A}$ risulta:
\[
    f(\A) = \sum_{n=0}^{\infty} \frac{\A^n}{n!} = \hat{Id} + \A + \frac{\A^2}{2} + \cdots
\]
\subsection*{Operatore Evoluzione Temporale}
Si indica con $\hat U(t,t_0)$ l'operatore di evoluzione temporale, che "fa avanzare" uno stato in un altro stato di un tempo $t-t_0$:
\[  
    |\psi,t\ket = \hat U(t,t_0) |\psi,t_0\ket
\]
L'operatore di evoluzione temporale deve rispettare alcune proprietà, come l'operatore di traslazione:
\begin{itemize}
    \item Deve essere lineare
    \item Deve essere unitario, per preservare la norma degli stati:
    \[
        \hat U^\dag (t,t_0) \hat U(t,t_0) = \hat{Id}
    \]
    \item Deve rispettare la proprietà di composizione:
    \[
        \hat U(t_2,t_0) = \hat U(t_2,t_1) \hat U(t_1,t_0)
    \]
    \item Deve essere continuo rispetto al tempo:
    \[
        \lim_{t \to t_0} \hat U(t,t_0) = \hat{Id}
    \]
\end{itemize}
Si consideri ora l'operatore infinitesimo di evoluzione temporale:
\[
    \hat U(t_0 + dt, t_0) = \hat{Id} - \frac i \hbar \hat H dt
\]
dove $\hat H$ è un operatore \textbf{indipendente dal tempo}, ed hermitiano così da garantire l'unitarietà di $\hat U$, la cui unità di misura è il Joule. \\
Ora espandiamo questo caso particolare:
\begin{gather*}
    \hat U(t + dt, t) \hat U(t, t_0) = \hat U(t + dt, t_0) \\
    (\hat{Id} - i / \hbar \hat H dt) \hat U(t, t_0) = \hat U(t + dt, t_0) \\
    \hat U(t + dt, t_0) - \hat U(t, t_0) = \frac{i}{\hbar} \hat H \hat U(t, t_0) dt
\end{gather*}
Che è molto simiele alla definizione di derivata; tuttavia siccome si tratta di operatori è necessario definire il comportamento della derivata in questo caso:
\[
    \deriv[t]{} \hat U(t, t_0) = \lim_{dt \to 0} \frac{\hat U(t + dt, t_0) - \hat U(t, t_0)}{dt} = -\frac{i}{\hbar} \hat H \hat U(t, t_0)
\]
che, posto $\hat H$ l'hamiltoniana solita, corrisponde alla dichicarazione dell'equazione di Schrödinger, infatti:
\begin{gather*}
    i \h \deriv[t]{} \hat U(t, t_0) |\psi\ket = \hat H \hat U(t, t_0) |\psi\ket \\
    i \h \deriv[t]{} \Ket {\psi(t)} = \hat H \Ket {\psi(t)}
\end{gather*}

Ora proviamo ad espandere l'operatore hamiltoniano e a porci nella base delle posizioni:
\begin{gather*}
    i \h \tderiv{} {\psi(\vec x,t)} = \Bra{\vec x}\frac{\hat p ^2}{2m} \Ket {\psi(t)} + \Bra{\vec x}\V \Ket {\psi(t)}
\end{gather*}
Se poi si ha che $\V$ è un potenziale locale (ovvero che dipende solo dalla posizione, e ne conserva gli autostati), risulta ancora:
\begin{gather*}
    \Bra{\vec x}\V \Ket {\psi(t)} = \rnint \Bra{\vec x}\V \ketbrabra{\vec x'} \Ket {\psi(t)} d^nx'k = \V(\vec x) \psi(\vec x, t)
\end{gather*}

Ora, dimostrato che l'hamiltoniana è proprio l'operatore che serve per ricostruire l'operatore di evoluzione temporale per tempi finiti, si può procedere come per le traslazioni spaziali, ripetendo a cascata $N$ volte l'operatore infinitesimo:
\begin{align*}
    \hat U(t, t_0) &= \lim_{N \to \infty} \left[ \hat U\left(t_0 + \frac{t - t_0}{N}, t_0\right) \right]^N \\
    &= \lim_{N \to \infty} \left[ \hat{Id} - \frac{i}{\hbar} \hat H \frac{t - t_0}{N} \right]^N \\
    &= e^{-\frac{i}{\hbar} \hat H (t - t_0)}
\end{align*}

Detto quetso non abbiamo mai ipotizzato che $\H$ potesse cambiare nel tempo: infatti, in questa forma, la definizione è valida solo se l'hamiltoniana è indipendente dal tempo.

\subsubsection*{Stati nella Base delle Energie}
Inoltre accade che se si sceglie come base per lo spazio di Hilbert gli autostati dell'hamiltoniano $\{\Ket{\vphi_j}, j \in J\}$, oltre a molte altre comodità che salteranno fuoir più avanti, l'operatore di evoluzione temporale risulta diagonale, in quanto è definito come somme di potenze di $\H$, che conservando tutti la stessa base:
\[
    \hat U(t, t_0) = \sum_j e^{-\frac{i}{\hbar} a_j (t - t_0)} |\vphi_j \ket \bra \vphi_j |
\]
e l'espressione dello stato $\psi$ in un tempo $t$ risulta:
\[
    |\psi, t\ket = \hat U(t, t_0) |\psi, t_0\ket = \sum_j e^{-\frac{i}{\hbar} a_j (t - t_0)} |\vphi_j \ket \bra \vphi_j | \psi, t_0\ket
\]
ma siccome gli stati $|\vphi_j \ket$ sono autostati dell'hamiltoniano, si ha che:
\begin{equation*}
    \hat H |\vphi_j \ket = a_j |\vphi_j \ket
    \implies |\vphi_j, t\ket = e^{-\frac{i}{\hbar} a_j (t - t_0)} |\vphi_j, t_0\ket
\end{equation*}
Al chè:
\[
    |\psi, t\ket = \sum_j c_j e^{-\frac{i}{\hbar} a_j (t - t_0)} |\vphi_j \ket
\]
Quest'equazione riassume tutta la dinamica quantistica, che è molto comodo, infatti:
\begin{align*}
    c_j(t) &= c_j(t_0) e^{-\frac{i}{\hbar} a_j (t - t_0)} \\
    |c_j(t)|^2 &= |c_j(t_0)|^2|e^{-\frac{i}{\hbar} a_j (t - t_0)}|^2 = |c_j(t_0)|^2
\end{align*}
ovvero le probabilità di trovare il sistema in uno stato autostato dell'hamiltoniano sono costanti nel tempo!
\subsection*{Valori medi degli osservabili nel tempo}
Dato uno stato $|\psi, t\ket$ che evolve nel tempo secondo l'operatore di evoluzione temporale, si può calcolare il valore medio di un osservabile $\hat A$ nel tempo $t$ come:
\begin{align*}
    \bra \hat A(t) \ket_ \psi &= \bra \psi, t | \hat A | \psi, t \ket \\
    &= \bra \psi, t_0 | \hat U^\dag(t, t_0) \hat A \hat U(t, t_0) | \psi, t_0 \ket \\
    &= \bra \psi, t_0 | e^{\frac{i}{\hbar} a_j (t - t_0)} \hat A e^{-\frac{i}{\hbar} a_j (t - t_0)} | \psi, t_0 \ket
\end{align*}
Ora si possono analizzare due casi:
\begin{itemize}
    \item $\hat A$ è costante del moto ($[\hat A, \hat H] = 0$). Banalmente risulta che:
    \[
        \bra \hat A(t)\ket_\psi = \bra \psi, t_0 | \U^\dag \hat A | \U \psi, t_0 \ket = \bra \psi, t_0 | \hat A | \psi, t_0 \ket
    \]
    in quanto si possono semplificare i due esponenziali, scambiando di posto $\A$ e $\U$:
    \[\comm{\A}{\H} = 0 \implies \comm{\A}{\H^n} = 0 \implies \comm{\A}{\U} = 0\]
    \item $\hat A$ non commuta con l'hamiltoniano. In questo caso si può inserire l'identità tra gli esponenziali e l'operatore:
    \[
        \bra \hat A(t) \ket_\psi = \bra \psi, t_0 | \hat A | \psi, t_0 \ket = \]\[
        = \sum_{i,j} \bra \psi, t_0|\psi_i\ket e^{\frac{i}{\hbar} (a_j - a_i) (t - t_0)} \bra \psi_i | \hat A | \vphi_j\ket \bra \vphi_j | \psi, t_0\ket = \]\[
        = \sum_{i,j} c_i^* c_j e^{\frac{i}{\hbar} (a_j - a_i) (t - t_0)} A_{ij}
    \]
    dove $A_{ij}$ è la matrice dell'operatore $\hat A$ nella base degli autostati dell'hamiltoniano. Non è un conto semplice, quindi ci si limita ad analizzare un paio di casi specifici:

    \textit{Si consideri uno stato del tipo $| \psi \ket = |\psi_0\ket$. Risulta facilmente che $\bra A(t)\ket_\psi = \bra A(t_0)\ket_\psi$.}

    \textit{Si consideri ora uno stato del tipo $| \psi \ket = \frac{1}{\sqrt{2}}|\psi_0\ket + \frac{1}{\sqrt{2}}|\psi_1\ket$. Stavolta risulta: (esercizioo) che differisce dalla precedente per un termine di interferenza. (si chiama $\omega_{ij}$ la differenza di energie degli autostati (gli autovalori) diviso $\hslash$)}
\end{itemize}
Questi esempi mettono in luce il fatto che se uno stato di partenza è una sovrapposizione di autostati, allora l'evoluzione nel tempo dipende solo dagli autostati considerati, e non dagli autostati che non erano considerati.




\section{Rappresentazioni alternative}
La rappresentazione di Schrödinger è quella che abbiamo usato fin'ora, in cui gli stati evolvono nel tempo secondo l'equazione di Schrödinger e gli operatori sono costanti nel tempo. (\textit{Schrödinger picture}).

È possibile introdurre anche un'altra rappresentazione, detta \textit{Heisenberg picture}\index{Rappresentazione di Heisenberg}, in cui sono gli operatori a evolvere nel tempo e gli stati sono costanti. \\
Si definisce l'operatore di Heisenberg come:
\begin{align*}
    \hat A_H(t) &= \hat U^\dag(t, t_0) \hat A \hat U(t, t_0) = \\
    &= e^{\frac{i}{\hbar} \hat H (t - t_0)} \hat A e^{-\frac{i}{\hbar} \hat H (t - t_0)}
\end{align*}
In cui il pedice $H$ indica che si tratta di un operatore in rappresentazione di Heisenberg. \\
Si può calcolare la derivata temporale di questo operatore:
\begin{align*}
    \deriv[t]{} \hat A_H(t) &= \deriv[t]{} \left( e^{\frac{i}{\hbar} \hat H (t - t_0)} \hat A e^{-\frac{i}{\hbar} \hat H (t - t_0)} \right) \\
    &= -\frac{i}{\hbar} \hat H e^{\frac{i}{\hbar} \hat H (t - t_0)} \hat A e^{-\frac{i}{\hbar} \hat H (t - t_0)} + e^{\frac{i}{\hbar} \hat H (t - t_0)} \hat A \left( -\frac{i}{\hbar} \hat H e^{-\frac{i}{\hbar} \hat H (t - t_0)} \right) \\
    &= -\frac{i}{\hbar} \left[ \hat H, \hat A_H(t) \right]
\end{align*}
Da cui l'equazione di Heisenberg:
\[
    i \hbar \deriv[t]{} \hat A_H(t) = \left[ \hat H, \hat A_H(t) \right]
\]
(ricavabile per esercizio).

\subsection*{Rappresentazione di Heisenberg: un esempio}
Si consideri l'operatore posizione e l'operatore momento di una particella libera in 1D:
\begin{align*}
    \left[ \hat x, \hat p \right] &= i \hbar \\
    \hat H &= \frac{\hat p^2}{2m}
\end{align*}
Ora si cerca
\[
    \left[ \hat x, \hat p^2 \right] = \hat x \hat p^2 - \hat p^2 \hat x = \hat p \left[\hat x, \hat p \right] + \left[ \hat x, \hat p \right] \hat p = 2 i \hbar \hat p
\]
Ora
\[
    \left[ \hat x, \hat p^3 \right] = \hat x \hat p^3 - \hat p^3 \hat x = \hat p^2 \left[\hat x, \hat p \right] + \left[ \hat x, \hat p \right] \hat p^2 = 3 i \hbar \hat p^2
\]
E in generale
\[
    \left[ \hat x, \hat p^n \right] = n i \hbar \hat p^{n-1}
\]
(si può dimostrare per induzione come esercizio). Questa operatzione è molto simile alla derivata, infatti si può scrivere:
\[
    \left[ \hat x, f(\hat p) \right] = i \hbar \deriv[\hat p]{} f(\hat p)
\]
E lo stesso vale scambiando $\hat p$ con $\hat x$. Non necessita dimostrazione il fatto che
\[
    \left[\hat x, f(\hat p)\right] = i\hbar \deriv[\hat p]{} f(\hat p);\ \ \left[\hat p, f(\hat x)\right] = -i\hbar \deriv[\hat x]{} f(\hat x)
\]